\chapter{Complexity — DAG Fanin}
%%% generated by the blueprint pipeline from TCSlib.Complexity.CircuitComplexity.DAGFanin
\section{Overview}

Declarations of the Lean module \texttt{TCSlib.Complexity.CircuitComplexity.DAGFanin}.

\section{Declarations}

\begin{definition}[Boolean operation of a gate]
\lean{BoolCircuit.GateKind.op}
\label{BoolCircuit.GateKind.op}
\leanok
\uses{BoolCircuit.GateKind}
Associate to each AND gate the Boolean conjunction of its two inputs, to each OR gate
their disjunction, and to each NOT gate the negation of its first input, ignoring the
second input in the NOT case.
\end{definition}

\begin{theorem}[Evaluation of concatenated gate inputs]
\lean{BoolCircuit.DAGGate.eval_append}
\label{BoolCircuit.DAGGate.eval_append}
\leanok
\uses{BoolCircuit.DAGGate, BoolCircuit.DAGGate.eval, BoolCircuit.GateKind, BoolCircuit.GateKind.op}
\difficulty{4}
Let $k$ be an AND or OR gate, let $l_1$ and $l_2$ be finite lists of natural-number
indices, and let $vals$ be a finite list of Boolean values. The gate of kind $k$ with
input list $l_1$ followed by $l_2$ evaluates to the result of applying conjunction or
disjunction, according to $k$, to the evaluations of the gates of kind $k$ with input
lists $l_1$ and $l_2$, respectively.
\end{theorem}

\begin{theorem}[Evaluation of a two-input gate]
\lean{BoolCircuit.DAGGate.eval_pair}
\label{BoolCircuit.DAGGate.eval_pair}
\leanok
\uses{BoolCircuit.DAGGate, BoolCircuit.DAGGate.eval, BoolCircuit.GateKind, BoolCircuit.GateKind.op}
\difficulty{4}
Let $k$ be a gate kind other than negation, let $a$ and $b$ be natural-number input
indices, and let $v$ be a list of Boolean values. The gate of kind $k$ with inputs $a,b$
evaluates to the binary operation associated with $k$ applied to the values at those
indices, taking an out-of-range index to have value false.
\end{theorem}

\begin{theorem}[Evaluation of a one-input gate]
\lean{BoolCircuit.DAGGate.eval_singleton}
\label{BoolCircuit.DAGGate.eval_singleton}
\leanok
\uses{BoolCircuit.DAGGate, BoolCircuit.DAGGate.eval, BoolCircuit.GateKind}
\difficulty{4}
Let $k$ be a gate kind other than negation, let $a$ be a natural-number input index, and
let $v$ be a list of Boolean values. The gate of kind $k$ with sole input $a$ evaluates
to the value of $v$ at index $a$, taking an out-of-range index to have value false.
\end{theorem}

\begin{theorem}[Evaluation is invariant under deduplication]
\lean{BoolCircuit.DAGGate.eval_dedup}
\label{BoolCircuit.DAGGate.eval_dedup}
\leanok
\uses{BoolCircuit.DAGGate, BoolCircuit.DAGGate.eval, BoolCircuit.GateKind}
\difficulty{4}
Let $k$ be a conjunction or disjunction gate, let $l$ be a finite list of natural-number
input indices, and let $vals$ be a finite list of Boolean values. The gate with kind $k$
and input list obtained from $l$ by removing repeated entries has the same evaluation on
$vals$ as the gate with input list $l$.
\end{theorem}

\begin{theorem}[Bound on a list maximum]
\lean{BoolCircuit.foldr_max_le₂}
\leanok
\difficulty{4}
Let $l$ be a finite list of natural numbers, let $f:br	obr$ be a function on the
natural numbers, and let $B$ be a natural number. If $f(a)\le B$ for every $a$ in $l$,
then the maximum of the values $f(a)$ over $a$ in $l$, taken to be $0$ when $l$ is
empty, is at most $B$.
\end{theorem}

\begin{theorem}[Bound by the maximum of mapped list values]
\lean{BoolCircuit.le_foldr_max₂}
\leanok
\difficulty{4}
Let $l$ be a list of natural numbers, let $f:\mathbb{N}\to\mathbb{N}$, and let $a\in l$.
Then $f(a)$ is at most the maximum of the values of $f$ on the entries of $l$, with the
maximum taken together with $0$.
\end{theorem}

\begin{definition}[Binary tree gate cost]
\lean{BoolCircuit.treeCost}
\label{BoolCircuit.treeCost}
\leanok
For a nonnegative integer $k$, define its tree cost to be $1$ if $k=0$, and $k-1$
otherwise.
\end{definition}

\begin{definition}[Recursive gate tree emission]
\lean{BoolCircuit.emitTree}
\label{BoolCircuit.emitTree}
\leanok
\uses{BoolCircuit.DAGGate, BoolCircuit.GateKind}
For a natural number $n$, a gate kind $k$, a list $v$ of natural-number vertex indices,
and a list $g$ of gates, define the emitted gate list and output vertex recursively. If
$v$ is empty, append a gate of kind $k$ with no arguments to $g$ and return vertex
$n+|g|$; if $v$ is a singleton, return $g$ and its sole vertex. Otherwise split $v$ into
its first $\lceil |v|/2\rceil$ entries and the remaining entries, recursively emit each
part in sequence, then append a gate of kind $k$ whose arguments are the two returned
vertices with duplicates removed, and return vertex $n$ plus the length of the resulting
gate list.
\end{definition}

\begin{definition}[Specification for a gate tree]
\lean{BoolCircuit.TreeSpec}
\label{BoolCircuit.TreeSpec}
\leanok
\uses{BoolCircuit.DAGGate, BoolCircuit.DAGGate.eval, BoolCircuit.GateKind, BoolCircuit.GatesAcyclic, BoolCircuit.runWith, BoolCircuit.treeCost, BoolCircuit.vertexDepth, BoolCircuit.vertexValue}
For an input count $n$, gate kind $k$, vertex list $v_s$, existing gate list $g_s$, and
proposed result $r=(r_g,r_v)$, the tree specification requires that $r_g$ extend $g_s$
by appending gates, that its gates be acyclic, and that $r_v<n+|r_g|$. For every Boolean
input assignment, the value at $r_v$ must equal the value of a gate of kind $k$ on $v_s$
evaluated using the original gates and input assignment. The result list has length at
most $|g_s|$ plus the tree cost of $|v_s|$, and the depth at $r_v$ is at most
$\lceil\log_2 |v_s|\rceil+1$ plus the maximum depth in the original circuit among
vertices in $v_s$ (with maximum taken together with $0$). Every appended gate has kind
$k$, has distinct arguments, and has at most two arguments.
\end{definition}

\begin{theorem}[Correctness and bounds for tree emission]
\lean{BoolCircuit.emitTree_spec}
\label{BoolCircuit.emitTree_spec}
\leanok
\uses{BoolCircuit.DAGGate, BoolCircuit.DAGGate.depth, BoolCircuit.DAGGate.eval, BoolCircuit.DAGGate.eval_append, BoolCircuit.DAGGate.eval_congr, BoolCircuit.DAGGate.eval_dedup, BoolCircuit.DAGGate.eval_pair, BoolCircuit.DAGGate.eval_singleton, BoolCircuit.GateKind, BoolCircuit.GatesAcyclic, BoolCircuit.GatesAcyclic.snoc, BoolCircuit.TreeSpec, BoolCircuit.emitTree, BoolCircuit.foldr_max_le₂, BoolCircuit.le_foldr_max₂, BoolCircuit.runWith, BoolCircuit.treeCost, BoolCircuit.vertexDepth, BoolCircuit.vertexDepth_append, BoolCircuit.vertexDepth_last, BoolCircuit.vertexValue, BoolCircuit.vertexValue_append, BoolCircuit.vertexValue_last}
\difficulty{7}
Let $n$ be a natural number, let $k$ be an AND or OR gate kind, and let $v_1,\ldots,v_m$
be a list of vertex indices. Let $G$ be an acyclic list of Boolean gates, and suppose
each $v_j<n+|G|$. The procedure that emits a binary tree of gates of kind $k$ returns an
extended acyclic gate list and a root index less than $n$ plus the length of that list.
On every Boolean input to the original $n$ vertices, the root value equals the value of
a single gate of kind $k$ applied to the listed vertices' values in $G$. The resulting
gate list has length at most $|G|+c(m)$, where $c(0)=1$ and $c(m)=m-1$ for $m>0$; the
root depth is at most $\lceil\log_2 m\rceil+1$ plus the maximum depth of the listed
vertices in $G$, with that maximum taken as $0$ for an empty list. Every appended gate
has kind $k$, has distinct input indices, and reads at most two vertices.
\end{theorem}

\begin{definition}[Binarization of a gate]
\lean{BoolCircuit.binarizeGadget}
\label{BoolCircuit.binarizeGadget}
\leanok
\uses{BoolCircuit.DAGGate, BoolCircuit.emitTree}
Given an input count $n$, a gate $g$, and an existing gate list $g_s$, if $g$ is a NOT
gate, append it unchanged and return vertex $n+|g_s|$. For an AND or OR gate,
recursively emit a tree of that kind on the distinct arguments of $g$, starting from
$g_s$.
\end{definition}

\begin{theorem}[Upper bound for tree cost]
\lean{BoolCircuit.treeCost_le}
\label{BoolCircuit.treeCost_le}
\leanok
\uses{BoolCircuit.treeCost}
\difficulty{2}
For every natural number $k$, the tree cost of $k$ is at most $k+2$.
\end{theorem}

\begin{theorem}[Correctness of an AND or OR gadget]
\lean{BoolCircuit.binarizeGadget_andor}
\label{BoolCircuit.binarizeGadget_andor}
\leanok
\uses{BoolCircuit.DAGGate, BoolCircuit.DAGGate.FaninTwo, BoolCircuit.DAGGate.eval_dedup, BoolCircuit.GadgetSpec, BoolCircuit.GateKind, BoolCircuit.GatesAcyclic, BoolCircuit.binarizeGadget, BoolCircuit.emitTree, BoolCircuit.emitTree_spec, BoolCircuit.foldr_max_le₂, BoolCircuit.le_foldr_max₂, BoolCircuit.treeCost_le, BoolCircuit.vertexDepth}
\difficulty{5}
Let $n,K$ be natural numbers, let $k$ be an AND or OR gate kind, and let $g$ be the gate
with kind $k$ and argument list $args$. Suppose the existing gate list $gs$ is acyclic,
every argument index is less than $n+|gs|$, and $|args|\le K$. Applying the binarization
construction to $g$ and $gs$ yields a gate list and output vertex that extend $gs$ and
form an acyclic circuit; the output vertex is less than $n$ plus the resulting list
length and computes $g$ on the values produced by $gs$ for every Boolean input
assignment. The resulting list has length at most $|gs|+|args|+2$, the output depth is
at most $\lceil\log_2 K\rceil+1$ plus the maximum depth among the argument vertices in
$gs$ (with empty maximum $0$), and every added gate is well formed with at most two
arguments.
\end{theorem}

\begin{theorem}[Correctness of gate binarization]
\lean{BoolCircuit.binarizeGadget_correct}
\label{BoolCircuit.binarizeGadget_correct}
\leanok
\uses{BoolCircuit.DAGGate.FaninTwo, BoolCircuit.DAGGate.depth, BoolCircuit.GadgetCorrect, BoolCircuit.GatesAcyclic.snoc, BoolCircuit.binarizeGadget, BoolCircuit.binarizeGadget_andor, BoolCircuit.runWith, BoolCircuit.vertexDepth, BoolCircuit.vertexDepth_last, BoolCircuit.vertexValue_last}
\difficulty{4}
For every natural number $K$, every gate kind $k$, argument list of length $\ell$, and
acyclic existing gate list $gs$ relative to $n$, assume each argument index is less than
$n+|gs|$, that $k=\mathrm{NOT}$ implies $\ell=1$, and that $\ell\le K$. Binarizing the
gate produces a gate list and output vertex extending $gs$ and forming an acyclic
circuit; the output vertex is less than $n$ plus the resulting list length and computes
the original gate's value on the values produced by $gs$ for every Boolean input
assignment. The resulting list has length at most $|gs|+\ell+2$, its output depth is at
most $\lceil\log_2 K\rceil+1$ plus the maximum depth among the argument vertices in $gs$
(with empty maximum $0$), and every added gate is well formed with at most two
arguments.
\end{theorem}

\begin{theorem}[Gate argument count bounded by circuit size]
\lean{BoolCircuit.DAGCircuit.args_length_le_size}
\label{BoolCircuit.DAGCircuit.args_length_le_size}
\leanok
\uses{BoolCircuit.DAGCircuit, BoolCircuit.DAGCircuit.IsWellFormed, BoolCircuit.DAGCircuit.size}
\difficulty{4}
Let $C$ be a finite Boolean circuit in which every gate has distinct arguments, and
every NOT gate has exactly one argument. Then each gate has at most $	ext{size}(C)$
arguments.
\end{theorem}

\begin{definition}[Circuit binarization]
\lean{BoolCircuit.DAGCircuit.binarize}
\label{BoolCircuit.DAGCircuit.binarize}
\leanok
\uses{BoolCircuit.DAGCircuit, BoolCircuit.DAGCircuit.IsWellFormed, BoolCircuit.DAGCircuit.args_length_le_size, BoolCircuit.DAGCircuit.rewrite, BoolCircuit.DAGCircuit.size, BoolCircuit.binarizeGadget, BoolCircuit.binarizeGadget_correct}
A well-formed circuit is one in which every gate has distinct arguments and every NOT
gate has exactly one argument. For any well-formed circuit, its binarization is the
circuit obtained by rewriting each gate using the gate-binarization construction, with
the resulting circuit retaining the same number of inputs.
\end{definition}

\begin{theorem}[Evaluation of the binarized circuit]
\lean{BoolCircuit.DAGCircuit.binarize_eval}
\label{BoolCircuit.DAGCircuit.binarize_eval}
\leanok
\uses{BoolCircuit.DAGCircuit.binarize, BoolCircuit.DAGCircuit.eval, BoolCircuit.DAGCircuit.rewrite_eval}
\difficulty{1}
Let $C$ be a directed acyclic Boolean circuit with $n$ input vertices, and suppose each
gate has pairwise distinct arguments and every negation gate has exactly one argument.
For every assignment $x$ of Boolean values to the inputs, the circuit obtained by
binarizing $C$ has the same output as $C$ on $x$.
\end{theorem}

\begin{theorem}[Fan-in of the binarized circuit]
\lean{BoolCircuit.DAGCircuit.binarize_isFaninTwo}
\label{BoolCircuit.DAGCircuit.binarize_isFaninTwo}
\leanok
\uses{BoolCircuit.DAGCircuit.IsFaninTwo, BoolCircuit.DAGCircuit.binarize, BoolCircuit.DAGCircuit.rewrite_new_gates}
\difficulty{1}
Let $C$ be a directed acyclic Boolean circuit with $n$ input vertices, and suppose each
gate has pairwise distinct arguments and every negation gate has exactly one argument.
After binarization, every gate still has pairwise distinct arguments, every negation
gate has exactly one argument, and every gate has at most two arguments.
\end{theorem}

\begin{theorem}[Size of the binarized circuit]
\lean{BoolCircuit.DAGCircuit.binarize_size_le}
\label{BoolCircuit.DAGCircuit.binarize_size_le}
\leanok
\uses{BoolCircuit.DAGCircuit.args_length_le_size, BoolCircuit.DAGCircuit.binarize, BoolCircuit.DAGCircuit.rewrite_length_le, BoolCircuit.DAGCircuit.size, BoolCircuit.DAGCircuit.sum_cost_le, BoolCircuit.binarizeGadget_correct}
\difficulty{4}
Let $C$ be a directed acyclic Boolean circuit with $n$ input vertices, and suppose each
gate has pairwise distinct arguments and every negation gate has exactly one argument.
The size of its binarization is at most $n+|G|(|C|+2)$, where $G$ is the list of gates
of $C$ and $|C|=n+|G|$.
\end{theorem}

\begin{theorem}[Depth of the binarized circuit]
\lean{BoolCircuit.DAGCircuit.binarize_depth_le}
\label{BoolCircuit.DAGCircuit.binarize_depth_le}
\leanok
\uses{BoolCircuit.DAGCircuit.binarize, BoolCircuit.DAGCircuit.depth, BoolCircuit.DAGCircuit.rewrite_depth_le, BoolCircuit.DAGCircuit.size}
\difficulty{1}
Let $C$ be a directed acyclic Boolean circuit with $n$ input vertices, and suppose each
gate has pairwise distinct arguments and every negation gate has exactly one argument.
If circuit depth is the depth of its output vertex, then binarizing $C$ gives a circuit
of depth at most $(\lceil\log_2(\mathrm{size}(C))\rceil+1)$ times the depth of $C$.
\end{theorem}

\begin{definition}[Emitting negation gates]
\lean{BoolCircuit.emitNots}
\label{BoolCircuit.emitNots}
\leanok
\uses{BoolCircuit.DAGGate}
Given a natural number $n$, a list of vertex indices, and a list of gates, recursively
append one NOT gate for each listed vertex, in list order. Return the resulting gate
list together with the indices assigned to those new gates, where each index is $n$ plus
the gate-list length before that gate is appended.
\end{definition}

\begin{definition}[Specification for emitted negations]
\lean{BoolCircuit.NotsSpec}
\label{BoolCircuit.NotsSpec}
\leanok
\uses{BoolCircuit.DAGGate, BoolCircuit.DAGGate.WellFormed, BoolCircuit.GatesAcyclic, BoolCircuit.vertexDepth, BoolCircuit.vertexValue}
Let $n$ be a natural number, let $a_1,\ldots,a_k$ be a list of vertex indices, let $G$
be a list of gates, and let $(G',v_1,\ldots,v_m)$ denote a pair consisting of a gate
list and a list of vertex indices. The specification requires that $G'$ extend $G$ by
appending gates, be acyclic, and have length $|G|+k$; the index list has length $k$,
contains distinct vertices, and each of its entries is at least $n+|G|$ and less than
$n+|G'|$. For every Boolean assignment to the $n$ input vertices, the values at these
indices are the negations of the values at $a_1,\ldots,a_k$, respectively; each such
vertex has depth at most one plus the maximum depth among those listed vertices, and
every appended gate is well formed and is a NOT gate.
\end{definition}

\begin{theorem}[Correctness of negation-gate emission]
\lean{BoolCircuit.emitNots_spec}
\label{BoolCircuit.emitNots_spec}
\leanok
\uses{BoolCircuit.DAGGate, BoolCircuit.DAGGate.depth, BoolCircuit.DAGGate.eval, BoolCircuit.GatesAcyclic, BoolCircuit.GatesAcyclic.snoc, BoolCircuit.NotsSpec, BoolCircuit.emitNots, BoolCircuit.foldr_max_le₂, BoolCircuit.le_foldr_max₂, BoolCircuit.vertexDepth, BoolCircuit.vertexDepth_append, BoolCircuit.vertexDepth_last, BoolCircuit.vertexValue, BoolCircuit.vertexValue_append, BoolCircuit.vertexValue_last}
\difficulty{6}
Let $n$ be a natural number, let $a_1,\ldots,a_m$ be a list of vertex indices, and
let $G$ be a list of Boolean gates such that every gate reads only earlier vertices:
gate $i$ reads indices less than $n+i$. Suppose every $a_j<n+|G|$. Then the procedure
that appends one negation gate for each listed index returns a gate list and
corresponding output indices with the following properties: the new gate list extends
$G$, is acyclic, and has length $|G|+m$; there are $m$ distinct output indices, each
between $n+|G|$ and strictly less than $n$ plus the new gate-list length. On every
Boolean input to the original $n$ vertices, these outputs are the negations of the
values at $a_1,\ldots,a_m$ in the original gate list. Each output has depth at most
one plus the maximum original depth of the listed vertices (with maximum taken as $0$
for an empty list), and every appended gate is a well-formed negation gate.
\end{theorem}

\begin{definition}[De Morgan gate gadget]
\lean{BoolCircuit.deMorganGadget}
\label{BoolCircuit.deMorganGadget}
\leanok
\uses{BoolCircuit.DAGGate, BoolCircuit.emitNots}
Given a natural number $n$, a gate $g$, and a gate list $G$, the gadget returns a gate
list and a vertex index. If $g$ is an OR gate, it first appends NOT gates for the
distinct vertices in the argument list, then an AND gate reading their new indices, then
a NOT gate reading the AND gate's index; it returns the resulting list and the index of
the final NOT gate. If $g$ is an AND gate, it appends an AND gate reading the distinct
arguments of $g$ and returns its index. If $g$ is a NOT gate, it appends $g$ itself and
returns its index.
\end{definition}

\begin{definition}[Well-formed gate without disjunction]
\lean{BoolCircuit.DAGGate.NoOr}
\label{BoolCircuit.DAGGate.NoOr}
\leanok
\uses{BoolCircuit.DAGGate, BoolCircuit.DAGGate.WellFormed}
For a natural number $K$ and a gate $g$, the condition holds exactly when $g$ is well
formed, is not an OR gate, and has at most $\max(1,K)$ arguments.
\end{definition}

\begin{theorem}[Correctness of the De Morgan gadget]
\lean{BoolCircuit.deMorganGadget_correct}
\label{BoolCircuit.deMorganGadget_correct}
\leanok
\uses{BoolCircuit.DAGGate, BoolCircuit.DAGGate.NoOr, BoolCircuit.DAGGate.depth, BoolCircuit.DAGGate.eval, BoolCircuit.DAGGate.eval_dedup, BoolCircuit.GadgetCorrect, BoolCircuit.GatesAcyclic, BoolCircuit.GatesAcyclic.snoc, BoolCircuit.deMorganGadget, BoolCircuit.emitNots, BoolCircuit.emitNots_spec, BoolCircuit.foldr_max_le₂, BoolCircuit.le_foldr_max₂, BoolCircuit.runWith, BoolCircuit.vertexDepth, BoolCircuit.vertexDepth_last, BoolCircuit.vertexValue, BoolCircuit.vertexValue_last}
\difficulty{6}
Fix natural numbers $n$ and $K$. For every gate $g$ and finite list $G$ of gates, assume
that each argument of the $i$th gate of $G$ is less than $n+i$, and each argument of $g$
is less than $n+|G|$. Assume also that if $g$ is a NOT gate then it has exactly one
argument, and that the number of arguments of $g$ is at most $K$. The De Morgan
construction applied to $g$ after $G$ produces a gate list extending $G$ and an output
vertex whose index is less than $n$ plus the length of the resulting list. The resulting
list is acyclic, and on every Boolean input the output vertex has the value of $g$
evaluated on the values obtained by running $G$ on that input. The resulting list has
length at most $|G|+|g	ext{'s argument list}|+2$, the output depth is at most three plus
the maximum depth of the vertices read by $g$ (taking this maximum to be zero when there
are no arguments), and every added gate is well formed, is not an OR gate, and has at
most $\max(1,K)$ arguments.
\end{theorem}

\begin{definition}[De Morgan transformation of a circuit]
\lean{BoolCircuit.DAGCircuit.deMorgan}
\label{BoolCircuit.DAGCircuit.deMorgan}
\leanok
\uses{BoolCircuit.DAGCircuit, BoolCircuit.DAGCircuit.IsWellFormed, BoolCircuit.DAGCircuit.args_length_le_size, BoolCircuit.DAGCircuit.rewrite, BoolCircuit.DAGCircuit.size, BoolCircuit.deMorganGadget, BoolCircuit.deMorganGadget_correct}
Let $C$ be a directed acyclic circuit with $n$ input vertices, and suppose each gate has
distinct arguments and every NOT gate has exactly one argument. The De Morgan
transformation of $C$ is the circuit obtained by rewriting its gates with the De Morgan
gate gadget.
\end{definition}

\begin{theorem}[Evaluation preserved by the De Morgan transformation]
\lean{BoolCircuit.DAGCircuit.deMorgan_eval}
\label{BoolCircuit.DAGCircuit.deMorgan_eval}
\leanok
\uses{BoolCircuit.DAGCircuit.deMorgan, BoolCircuit.DAGCircuit.eval, BoolCircuit.DAGCircuit.rewrite_eval}
\difficulty{1}
Let $C$ be a finite acyclic Boolean circuit with $n$ input vertices, and suppose every
gate has distinct arguments and every NOT gate has exactly one argument. For every
Boolean input assignment, the circuit obtained by the De Morgan transformation has the
same output as $C$.
\end{theorem}

\begin{theorem}[Well-formedness after the De Morgan transformation]
\lean{BoolCircuit.DAGCircuit.deMorgan_isWellFormed}
\label{BoolCircuit.DAGCircuit.deMorgan_isWellFormed}
\leanok
\uses{BoolCircuit.DAGCircuit.IsWellFormed, BoolCircuit.DAGCircuit.deMorgan, BoolCircuit.DAGCircuit.rewrite_new_gates}
\difficulty{1}
Let $C$ be a finite acyclic Boolean circuit with $n$ input vertices, and suppose every
gate has distinct arguments and every NOT gate has exactly one argument. Then every gate
in the circuit obtained by the De Morgan transformation also has distinct arguments, and
each of its NOT gates has exactly one argument.
\end{theorem}

\begin{theorem}[No OR gates after the De Morgan transformation]
\lean{BoolCircuit.DAGCircuit.deMorgan_kind_ne_or}
\label{BoolCircuit.DAGCircuit.deMorgan_kind_ne_or}
\leanok
\uses{BoolCircuit.DAGCircuit.deMorgan, BoolCircuit.DAGCircuit.rewrite_new_gates}
\difficulty{1}
Let $C$ be a finite acyclic Boolean circuit with $n$ input vertices, and suppose every
gate has distinct arguments and every NOT gate has exactly one argument. Then the
circuit obtained by the De Morgan transformation contains no OR gates.
\end{theorem}

\begin{theorem}[Size bound for the De Morgan transformation]
\lean{BoolCircuit.DAGCircuit.deMorgan_size_le}
\label{BoolCircuit.DAGCircuit.deMorgan_size_le}
\leanok
\uses{BoolCircuit.DAGCircuit.args_length_le_size, BoolCircuit.DAGCircuit.deMorgan, BoolCircuit.DAGCircuit.rewrite_length_le, BoolCircuit.DAGCircuit.size, BoolCircuit.DAGCircuit.sum_cost_le, BoolCircuit.deMorganGadget_correct}
\difficulty{4}
Let $C$ be a finite acyclic Boolean circuit with $n$ input vertices, and suppose every
gate has distinct arguments and every NOT gate has exactly one argument. If the size of
a circuit is its number of input vertices plus its number of gates, then the transformed
circuit has size at most $n + |G|(s+2)$, where $|G|$ is the number of gates of $C$ and
$s=n+|G|$ is its size.
\end{theorem}

\begin{theorem}[Depth of the De Morgan transformation]
\lean{BoolCircuit.DAGCircuit.deMorgan_depth_le}
\label{BoolCircuit.DAGCircuit.deMorgan_depth_le}
\leanok
\uses{BoolCircuit.DAGCircuit.deMorgan, BoolCircuit.DAGCircuit.depth, BoolCircuit.DAGCircuit.rewrite_depth_le}
\difficulty{1}
Let $C$ be a finite acyclic Boolean circuit with $n$ input vertices, and suppose every
gate has distinct arguments and every NOT gate has exactly one argument. If circuit
depth is the maximum number of gates along a path to the output, then the circuit
obtained by the De Morgan transformation has depth at most three times the depth of $C$.
\end{theorem}
